\section{Altitude Hold PID}

Tầng altitude điều khiển độ cao $z$ của drone. Khác với Position PID, tầng này
không tạo góc nghiêng mà tạo PWM base cho cả bốn motor.

\subsection{Sai số độ cao và vận tốc đứng}

Gọi độ cao mong muốn là:

\begin{equation}
    z_d
\end{equation}

Độ cao hiện tại là:

\begin{equation}
    z
\end{equation}

Sai số độ cao:

\begin{equation}
    e_z = z_d - z
\end{equation}

Gọi vận tốc đứng mong muốn là $v_{zd}$ và vận tốc đứng hiện tại là $v_z$.
Sai số vận tốc đứng:

\begin{equation}
    e_{v_z} = v_{zd} - v_z
\end{equation}

Trong project, $v_{zd}$ mặc định bằng $0$, nghĩa là khi giữ độ cao thì drone
không mong muốn tiếp tục đi lên hoặc đi xuống.

\subsection{Tích phân độ cao}

Thành phần tích phân của sai số độ cao:

\begin{equation}
    I_z(k) = I_z(k-1) + e_z(k)\Delta t
\end{equation}

Project giới hạn:

\begin{equation}
    -I_{z,max} \leq I_z \leq I_{z,max}
\end{equation}

Trong code:

\begin{equation}
    I_{z,max}=12.0
\end{equation}

Giới hạn này giúp tránh integral windup khi drone ở xa target hoặc khi PWM đã
bị kẹp ở biên.

\subsection{Output của Altitude PID}

Altitude PID tạo ra lượng thay đổi PWM:

\begin{equation}
    \Delta PWM_z
    =
    K_{p_z}e_z
    +
    K_{d_z}e_{v_z}
    +
    K_{i_z}I_z
\end{equation}

Base PWM được tính bằng:

\begin{equation}
    PWM_{base}
    =
    PWM_{hover}
    +
    \Delta PWM_z
\end{equation}

Sau đó được giới hạn trong khoảng:

\begin{equation}
    PWM_{min}
    \leq
    PWM_{base}
    \leq
    PWM_{max}
\end{equation}

Với:

\begin{equation}
    PWM_{min}=1100,\quad PWM_{max}=1940
\end{equation}

Do đó công thức đầy đủ trong code là:

\begin{equation}
    PWM_{base}
    =
    clamp(
    PWM_{hover}
    +
    K_{p_z}e_z
    +
    K_{d_z}e_{v_z}
    +
    K_{i_z}I_z,
    1100,
    1940
    )
\end{equation}

\subsection{Ý nghĩa vật lý}

$PWM_{base}$ là mức ga chung cho bốn motor. Nếu drone thấp hơn target thì
$e_z>0$, khi đó $\Delta PWM_z$ tăng và tổng lực nâng tăng. Nếu drone cao hơn
target thì $e_z<0$, PWM nền giảm để drone hạ xuống.

Tầng altitude chỉ lo tổng lực nâng, còn việc phân bố lực giữa các motor để giữ
roll, pitch, yaw sẽ do mixer xử lý.

\section{Attitude Hold PID}

Tầng attitude giữ roll và pitch của drone. Input của tầng này là góc mong muốn
từ position controller:

\begin{equation}
    \phi_d,\quad \theta_d
\end{equation}

và góc hiện tại:

\begin{equation}
    \phi,\quad \theta
\end{equation}

Trong đó $\phi$ là roll và $\theta$ là pitch.

\subsection{Sai số roll/pitch}

Sai số roll:

\begin{equation}
    e_\phi = wrap(\phi_d - \phi)
\end{equation}

Sai số pitch:

\begin{equation}
    e_\theta = wrap(\theta_d - \theta)
\end{equation}

Hàm $wrap$ giữ sai số góc trong khoảng $[-\pi,\pi]$.

\subsection{Thành phần vi phân dùng vận tốc góc}

Trong code, thành phần D của attitude PID không tính bằng:

\begin{equation}
    \frac{d e_\phi}{dt},\quad \frac{d e_\theta}{dt}
\end{equation}

mà dùng trực tiếp vận tốc góc body:

\begin{equation}
    p = \dot{\phi},\quad q = \dot{\theta}
\end{equation}

Nếu dùng PID chuẩn, thành phần D là:

\begin{equation}
    K_{d_\phi}\frac{d e_\phi}{dt}
    =
    K_{d_\phi}\left(\frac{d \phi_d}{dt} - \frac{d \phi}{dt}\right)
    =
    K_{d_\phi}(\dot{\phi}_d - \dot{\phi})
\end{equation}

Nếu roll target thay đổi chậm: $\dot{\phi}_d \approx 0$ và $p \approx \dot{\phi}$


Vì vậy:

\begin{equation}
    K_{d_\phi}\frac{d e_\phi}{dt}
    \approx
    -K_{d_\phi}\dot{\phi}_d
    =
    -K_{d_\phi}p
\end{equation}

Tương tự với pitch, đó là lí do công thức có dấu trừ:

\begin{equation}
    \Delta PWM_\phi
    =
    K_{p_\phi}e_\phi
    +
    K_{i_\phi}I_\phi
    -
    K_{d_\phi}p
\end{equation}

\begin{equation}
    \Delta PWM_\theta
    =
    K_{p_\theta}e_\theta
    +
    K_{i_\theta}I_\theta
    -
    K_{d_\theta}q
\end{equation}

Khi drone đang quay nhanh theo một chiều, thành phần $-K_d p$ hoặc $-K_d q$
tạo lực hãm để giảm dao động.

\subsection{Tích phân và giới hạn output}

Tích phân roll/pitch:

\begin{equation}
    I_\phi(k)=I_\phi(k-1)+e_\phi(k)\Delta t
\end{equation}

\begin{equation}
    I_\theta(k)=I_\theta(k-1)+e_\theta(k)\Delta t
\end{equation}

Trong code:

\begin{equation}
    -1.0 \leq I_\phi,I_\theta \leq 1.0
\end{equation}

Output PWM correction cũng được kẹp:

\begin{equation}
    -260
    \leq
    \Delta PWM_\phi,\Delta PWM_\theta
    \leq
    260
\end{equation}

Các giới hạn này làm cho attitude controller không thể yêu cầu mixer tạo ra
lệnh quá mạnh so với khả năng motor.

\section{Yaw Hold PID}

Yaw controller giữ hướng $\psi$ của drone. Trong project, yaw controller có thể
bật/tắt bằng biến \texttt{enabled}.

\subsection{Sai số yaw}

Yaw target:

\begin{equation}
    \psi_d
\end{equation}

Yaw hiện tại:

\begin{equation}
    \psi
\end{equation}

Sai số yaw:

\begin{equation}
    e_\psi = wrap(\psi_d-\psi)
\end{equation}

Việc wrap góc đặc biệt quan trọng với yaw vì yaw có thể đi qua biên
$-\pi$ và $\pi$.

\subsection{Công thức yaw correction}

Gọi yaw rate là:

\begin{equation}
    r = \dot{\psi}
\end{equation}

Tích phân yaw:

\begin{equation}
    I_\psi(k)=I_\psi(k-1)+e_\psi(k)\Delta t
\end{equation}

Output của yaw PID là PWM correction:

\begin{equation}
    \Delta PWM_\psi
    =
    K_{p_\psi}e_\psi
    +
    K_{i_\psi}I_\psi
    -
    K_{d_\psi}r
\end{equation}

Trong code, output này không phải torque trực tiếp. Nó là lượng PWM sẽ được
mixer cộng/trừ lên các motor để tạo chênh lệch torque phản lực.

Giới hạn tích phân:

\begin{equation}
    -0.8 \leq I_\psi \leq 0.8
\end{equation}

Giới hạn output:

\begin{equation}
    -120
    \leq
    \Delta PWM_\psi
    \leq
    120
\end{equation}

\subsection{Vai trò của yaw correction trong quadcopter}

Roll và pitch được tạo chủ yếu bằng chênh lệch lực nâng giữa các motor. Yaw
được tạo bằng chênh lệch torque phản lực của các cánh quạt quay ngược chiều
nhau. Vì vậy mixer sẽ cộng $\Delta PWM_\psi$ vào nhóm motor này và trừ khỏi
nhóm motor kia.

Nếu yaw hold tắt, project đặt:

\begin{equation}
    \Delta PWM_\psi = 0
\end{equation}

Khi đó mixer chỉ còn giữ độ cao, roll và pitch.
